
\chapter{\hard\ \trSerreDuality{}的证明}
\label{ch:serre-duality}

\epigraph{稍微琢磨了一下你的对偶定理，我发现它的普遍表述几乎是显然的，而且我刚验证了：（就射影空间而言）它已隐含在你用 Ext 给出 $\mathrm{T}^q(M)$ 的那条定理之中。（我真有种强烈的印象——混蛋——你第 3 章的 §3 和 §4 不用任何计算就能做出来。）}{——A.~\trGrothendieck，致 J.-P.~\trSerre 书信，1955 年 12 月 15 日 [GrS, p.~19]}

我们与\trSerreDuality{}初见于 \S18.5（甚至\mbox{定理~18.1.2}），此后已一再见到它的用处。我们将证明\mbox{定理~18.5.1}（\mbox{推论~29.3.10}，辅以\mbox{练习~29.4.K}；见\mbox{注记~29.4.10}），以及若干更强的版本，而在证明之后，我们还会心生证明更多东西的渴望。我们给出三种陈述（向量丛的\trSerreDuality；$\Hom$ 的\trSerreDuality；$\mathrm{Ext}$ 的\trSerreDuality），各有两个版本（函子版本与迹版本）。（这些名称是个人习惯用法，并非通行标准。）
之所以给出好几种变体，理由有二。其一，较容易的陈述会更容易证明，而最难的陈述我们在这里将无力证明。其二，它们或许有助于你积累这样的经验：知道应当期望什么，以及知道要尝试证明什么。


\section{期望清单}

本章自始至终在一个域 $k$ 上工作；$X$ 将是一个\trProjective{} $k$-\trScheme，而 $n=\dim X$。我们将会想要 $X$ 上的一个\trCoherentSheaf{} $\boxed{\omega}$（或者，说得更精确些，$\boxed{\omega_X}$；甚至，更好的是 $\omega_{X/k}$），也就是\textbf{\trCanonical{}层}或\textbf{\trDualizingSheaf}，它将在\trDuality{}的诸陈述中扮演角色。为了得到最好的那些陈述，我们还会想要一个\textbf{迹\trMorphism}
\begin{equation*}\tag{29.1.0.1}\label{eq:trace-morphism-hoped}
\boxed{t\colon\mathrm{H}^n(X,\omega_X)\longrightarrow k.}
\end{equation*}
我们最终将在 \S29.1.5 中定义 $\omega$ 与 $t$，但首先我们来讨论一下，我们希望它们具有哪些性质。

\textbf{29.1.1. \trDesideratum{}：对光滑簇而言，\trCotangentBundle{}的\trDeterminant{}是\trDualizingSheaf。}\footnote{译注：节题 Desiderata（期望清单）与本条目头 Desideratum（愿望）为同源词：本节罗列的正是我们期望对偶性应当具有的一系列性质，故条目均以「愿望」为头。——译者注}\quad 若 $X$ \trSmooth，则在此情形我们将会想要 $\boxed{\omega_X=\mathscr{K}_X}$（回想 $\mathscr{K}_X=\det\Omega_X$）——亦即「\trCanonical{}丛是 \trSerre{}对偶化的」这一奇迹（\S18.5.2）。特别地，$\omega_X$ 是一个\trInvertibleSheaf。这件事证明起来会令人不安地比我们所展示的那些基本\trDuality{}陈述更难；我们只能留待以后再处理它（\mbox{练习~29.4.K}）。但我们会证明更多的结果，比如：若 $X$ 是一个\trSmooth\trVariety{}中的\trRegularEmbedding，则 $\omega_X$ 是一个\trInvertibleSheaf（\mbox{练习~29.4.J}）。

\textbf{29.1.2. \trDesideratum{}：\trVectorBundle{}的\trSerreDuality。}\quad \trDuality{}的第一个版本——它与\mbox{\trDesideratum{}~29.1.1} 一起给出\mbox{定理~18.5.1}——如下：若 $\mathscr{F}$ 是有限秩的局部自由层，则我们有函子性的\trIsomorphism
\begin{equation*}\tag{29.1.2.1}\label{eq:vector-bundle-duality-iso}
\mathrm{H}^i(X,\mathscr{F})\cong\mathrm{H}^{n-i}(X,\mathscr{F}^{\vee}\otimes\omega_X)^{\vee}.
\end{equation*}
更精确地说，我们要构造的是\emph{某个特定的}\trIsomorphism{} \eqref{eq:vector-bundle-duality-iso}，或者等价地，某个特定的\trPerfectPairing{} $\mathrm{H}^i(X,\mathscr{F})\times\mathrm{H}^{n-i}(X,\mathscr{F}^{\vee}\otimes\omega_X)\to k$。这个\trIsomorphism{}将是随 $\mathscr{F}$ 函子的，也就是说，它给出协变函子之间的一个\trNaturalIsomorphism
\begin{equation*}\tag{29.1.2.2}\label{eq:vector-bundle-functorial-iso}
\boxed{\mathrm{H}^i(X,\cdot)\overset{\sim}{\longrightarrow}\mathrm{H}^{n-i}(X,\cdot^{\vee}\otimes\omega_X)^{\vee}.}
\end{equation*}
我们称之为\textbf{\trVectorBundle{}的函子版本\trSerreDuality}。

更好的是，\trCohomology{}中还应当有一个上积（cup product），可用它来构造一个映射 $\mathrm{H}^i(X,\mathscr{F})\times\mathrm{H}^{n-i}(X,\mathscr{F}^{\vee}\otimes\omega_X)\to\mathrm{H}^n(X,\omega_X)$。它应当随 $\mathscr{F}$ 函子（其含义是：我们得到函子之间的一个\trNaturalTransformation{} $\mathrm{H}^i(X,\cdot)\to\mathrm{H}^{n-i}(X,\cdot^{\vee}\otimes\omega_X)^{\vee}\otimes\mathrm{H}^n(X,\omega_X)$）。将它与\trTraceMap \eqref{eq:trace-morphism-hoped} 相结合，我们就得到一个映射
\begin{equation*}\tag{29.1.2.3}\label{eq:vector-bundle-trace-composition}
\begin{tikzcd}[row sep=3.2em, column sep=3.4em]
\mathrm{H}^i(X,\mathscr{F})\times\mathrm{H}^{n-i}(X,\mathscr{F}^{\vee}\otimes\omega_X) \ar[r] & \mathrm{H}^n(X,\mathscr{F}\otimes\mathscr{F}^{\vee}\otimes\omega_X) \ar[d] \\
& \mathrm{H}^n(X,\omega_X) \ar[r, "t"] & k,
\end{tikzcd}
\end{equation*}
它若是一个\trPerfectPairing，就应当给出 \eqref{eq:vector-bundle-functorial-iso}。若复合映射 \eqref{eq:vector-bundle-trace-composition} 是一个\trPerfectPairing，我们就说 $X$（连同附加数据 $(\omega_X,t)$）满足\textbf{\trVectorBundle{}\trSerreDuality{}的迹版本}。（在\mbox{定理~18.5.1} 的陈述之后，我们曾暗示过这一点。）

事实上，\trVectorBundle{}\trSerreDuality{}的迹版本（从而函子版本）在 $X$ 为 \trCM{}时成立，特别地，在 $X$ \trSmooth{}时成立。我们将证明函子版本（见\mbox{推论~29.3.10}）；至于迹版本可以如何证明，我们会在\mbox{注记~29.3.15} 中给出一点提示。

\textbf{29.1.3. \trDesideratum{}：更一般的 $X$ 与 $\mathscr{F}$ 的\trDuality。}\quad 一种较弱的\trDuality{}将在较弱的假设下成立。我们将证明：对纯 $n$ 维\trProjective{} $k$-\trScheme{} $X$ 上的任意\trCoherentSheaf{} $\mathscr{F}$（无需 \trCM{}假设），存在一个\emph{函子性的\trIsomorphism}（随 $\mathscr{F}$ 函子）
\begin{equation*}\tag{29.1.3.1}\label{eq:hom-functorial-iso}
\boxed{\mathrm{Hom}(\mathscr{F},\omega_X)\overset{\sim}{\longrightarrow}\mathrm{H}^n(X,\mathscr{F})^{\vee}.}
\end{equation*}
我们称之为\textbf{$\mathrm{Hom}$ 的函子版本\trSerreDuality}。

与\trVectorBundle{}的\trSerreDuality{}相平行，对于这个\trPerfectPairing，我们有一个自然的候选者：
\begin{equation*}\tag{29.1.3.2}\label{eq:hom-trace-pairing}
\boxed{\mathrm{Hom}(\mathscr{F},\omega_X)\times\mathrm{H}^n(X,\mathscr{F})\longrightarrow\mathrm{H}^n(X,\omega_X)\overset{t}{\longrightarrow}k,}
\end{equation*}
其中\trTraceMap $t$ 是 $\mathrm{H}^n(X,\omega_X)$ 上的某个线性泛函。若这个复合是一个\trPerfectPairing，我们就说 $X$（连同附加数据 $(\omega_X,t)$）满足\textbf{$\mathrm{Hom}$ 的\trSerreDuality{}的迹版本}。与\trVectorBundle{}的\trSerreDuality{}的迹版本不同，这里的「上积」映射 $\mathrm{Hom}(\mathscr{F},\omega_X)\times\mathrm{H}^n(X,\mathscr{F})\to\mathrm{H}^n(X,\omega_X)$ 我们已经知道是什么：$\mathrm{Hom}(\mathscr{F},\omega_X)$ 的一个元素 $[\sigma\colon\mathscr{F}\to\omega_X]$，借 $\mathrm{H}^n(X,\cdot)$ 的协变性（见 \S18.1），诱导出一个映射 $\mathrm{H}^n(X,\mathscr{F})\to\mathrm{H}^n(X,\omega_X)$。所得的配对显然随 $\mathscr{F}$ 函子，因此 $\mathrm{Hom}$ 的\trSerreDuality{}的迹版本蕴含函子版本。而与\trVectorBundle{}的情形不同的是，$\mathrm{Hom}$ 的函子版本\trSerreDuality{}反过来也蕴含迹版本：

\textbf{29.1.A.（练习）}\quad 证明：$\mathrm{Hom}$ 的\trSerreDuality{}的函子版本蕴含迹版本。换言之，\trTraceMap{}已经隐含在函子性的\trIsomorphism{} $\mathrm{Hom}(\mathscr{F},\omega_X)\overset{\sim}{\longrightarrow}\mathrm{H}^n(X,\mathscr{F})^{\vee}$ 之中。提示：考虑（由函子性得来的）交换图
\[
\begin{tikzcd}[row sep=3.2em, column sep=4em]
\mathrm{Hom}(\mathscr{F},\omega_X) \ar[r] & \mathrm{H}^n(X,\mathscr{F})^{\vee} \\
\mathrm{Hom}(\omega_X,\omega_X) \ar[u] \ar[r] & \mathrm{H}^n(X,\omega_X)^{\vee}. \ar[u]
\end{tikzcd}
\]

\textbf{29.1.4. 开始动手。}\quad 我们先来定义其中的若干术语。

\textbf{29.1.5.}\ \emph{\trhDefinition{}。}\quad 设 $X$ 是\trDimension{}为 $n$ 的\trProjective{} $k$-\trScheme。一个\trCoherentSheaf{} $\omega=\omega_X=\omega_{X/k}$ 连同一个映射 $t\colon\mathrm{H}^n(X,\omega_X)\to k$ 合起来称作\textbf{对偶化的}，若自然映射（对照 \eqref{eq:hom-functorial-iso}）
\begin{equation*}\tag{29.1.5.1}\label{eq:dualizing-pairing}
\mathrm{Hom}(\mathscr{F},\omega_X)\times\mathrm{H}^n(X,\mathscr{F})\longrightarrow\mathrm{H}^n(X,\omega_X)\overset{t}{\longrightarrow}k
\end{equation*}
对 $X$ 上的一切\trCoherentSheaf{} $\mathscr{F}$ 都是\trPerfectPairing。我们称 $\omega_X$ 为\textbf{\trCanonical{}层}，称 $t$ 为\textbf{\trTraceMap}。（早先关于 $\omega_X$ 与 $t$ 的讨论只是愿望式的；现在这则是一个定义。）若 $X$ 具有这样一个 $(\omega_X,t)$，我们就说 $X$ 满足\textbf{$\mathrm{Hom}$ 的\trSerreDuality}。

\textbf{29.1.6.}\ \emph{记号旁注。}\quad 此处定义的\trCanonical{}层常被称作\textbf{\trDualizingSheaf}。这个术语并不算糟，但最好的做法或许是：把这个称呼留给满足\trSerreDuality{}最好版本——即 $\mathrm{Ext}$ 的迹版本（\S29.1.13）——的层，而那已超出本书的范围。类似地，为了精确起见，\mbox{定义~29.1.5} 中的「对偶化的」或许更宜称作「弱对偶化的」。

此处的「\trCanonical{}层」$\omega_X$，与 \S21.5.3 中光滑簇的「\trCanonical{}丛」$\mathscr{K}_X=\det\Omega_X$，在术语上存在一个表面上的冲突。但 \S29.4 的要点正是要表明：在两者都适用的场合，这两个概念是同一个东西。

\textbf{29.1.7.}\quad 下述\trhProposition{}为「the canonical sheaf（这个\trCanonical{}层）」这一说法中使用定冠词「the」（而非「a」）提供了正当性。

\textbf{29.1.8. \trhProposition。}\ ——\ \emph{若 $X$ 存在\trCanonical{}层与迹 $(\omega_X,t)$，则这份数据在相差唯一\trIsomorphism{}的意义下唯一。}

\emph{证明。}\quad 设我们有两个这样的 $(\omega_X,t)$ 与 $(\omega'_X,t')$。由那两个\trMorphism
\begin{equation*}\tag{29.1.8.1}\label{eq:two-trace-data}
\begin{gathered}
\mathrm{Hom}(\mathscr{F},\omega_X)\times\mathrm{H}^n(X,\mathscr{F})\longrightarrow\mathrm{H}^n(X,\omega_X)\overset{t}{\longrightarrow}k,\\[1.4em]
\mathrm{Hom}(\mathscr{F},\omega'_X)\times\mathrm{H}^n(X,\mathscr{F})\longrightarrow\mathrm{H}^n(X,\omega'_X)\overset{t'}{\longrightarrow}k,
\end{gathered}
\end{equation*}
我们得到一个自然的\trBijection{} $\mathrm{Hom}(\mathscr{F},\omega_X)\cong\mathrm{Hom}(\mathscr{F},\omega'_X)$，它随 $\mathscr{F}$ 函子。由典型的\trUniversalProperty{}论证（\mbox{练习~1.2.Z}），这诱导出一个（唯一的）\trIsomorphism{} $\omega_X\overset{\sim}{\hookrightarrow}\omega'_X$。由 \eqref{eq:two-trace-data}，在这一同构之下，两个\trTraceMap $t$ 与 $t'$ 也必定相同。\hfill$\square$

我们将在\mbox{推论~29.3.12} 中证明 $\mathrm{Hom}$ 的\trSerreDuality{}的函子版本与迹版本。射影空间这一特殊情形将是关键的原料。我们现在就来证明这个情形。

\textbf{29.1.9.}\ \emph{$\mathbb{P}^n_k$ 上（$\mathrm{Hom}$ 的）\trSerreDuality。}\quad 定义 $\omega_{\mathbb{P}^n_k}$（为方便起见，或径记作 $\omega$）为 $\OO_{\mathbb{P}^n_k}(-n-1)$。令 $t$ 为任一\trIsomorphism{} $\mathrm{H}^n(\mathbb{P}^n_k,\omega_{\mathbb{P}^n_k})\overset{\sim}{\hookrightarrow}k$（\mbox{定理~18.1.2}(b)）。正如这个记号所暗示的，$(\omega_{\mathbb{P}^n_k},t)$ 对射影空间 $\mathbb{P}^n_k$ 而言将会是对偶化的。

\textbf{29.1.B.（练习）}\quad 设 $\mathscr{F}=\OO_{\mathbb{P}^n_k}(m)$。证明：自然映射 \eqref{eq:dualizing-pairing} 是一个\trPerfectPairing。（提示：亲自动手做。参看\mbox{定理~18.1.2} 之后的讨论。）进而证明：若 $\mathscr{F}$ 是 $\mathbb{P}^n_k$ 上若干\trLineBundle{}的一个\trDirectSum，则自然映射 \eqref{eq:dualizing-pairing} 是一个\trPerfectPairing。

\textbf{29.1.10. \trhProposition。}\ ——\ \emph{$\mathrm{Hom}$ 的\trSerreDuality{}的函子版本（从而由\mbox{练习~29.1.A}，迹版本亦然）对 $\mathbb{P}^n_k$ 成立。}

\emph{证明。}\quad 我们希望证明：由上积配对与\trIsomorphism{} $t$ 的选取而来的（函子性的）自然映射
\[
\mathrm{Hom}(\mathscr{F},\omega_{\mathbb{P}^n_k})\overset{\sim}{\longrightarrow}\mathrm{H}^n(\mathbb{P}^n_k,\mathscr{F})^{\vee},
\]
是一个\trIsomorphism（对照 \eqref{eq:hom-functorial-iso}）。取定 $\mathbb{P}^n_k$ 上的一个\trCoherentSheaf{} $\mathscr{F}$。依据\mbox{定理~16.1.1}，我们可以将 $\mathscr{F}$ 呈现为
\begin{equation*}\tag{29.1.10.1}\label{eq:presentation-of-f}
0\longrightarrow\mathscr{G}\longrightarrow\mathscr{E}\longrightarrow\mathscr{F}\longrightarrow 0,
\end{equation*}
其中 $\mathscr{E}$ 是有限多个\trLineBundle{}的一个\trDirectSum，而 $\mathscr{G}$ 凝聚。对 \eqref{eq:presentation-of-f} 施加\trLeftExact{}函子 $\mathrm{Hom}(\cdot,\omega_{\mathbb{P}^n_k})$，我们便有\trExactSequence
\begin{equation*}\tag{29.1.10.2}\label{eq:hom-left-exact}
0\longrightarrow\mathrm{Hom}(\mathscr{F},\omega_{\mathbb{P}^n_k})\longrightarrow\mathrm{Hom}(\mathscr{E},\omega_{\mathbb{P}^n_k})\longrightarrow\mathrm{Hom}(\mathscr{G},\omega_{\mathbb{P}^n_k}).
\end{equation*}
对 \eqref{eq:presentation-of-f} 取\trCohomology{}的\trLongExactSequence{}并对偶化，我们便有\trExactSequence
\begin{equation*}\tag{29.1.10.3}\label{eq:cohom-dual-exact}
0\longrightarrow\mathrm{H}^n(\mathbb{P}^n_k,\mathscr{F})^{\vee}\longrightarrow\mathrm{H}^n(\mathbb{P}^n_k,\mathscr{E})^{\vee}\longrightarrow\mathrm{H}^n(\mathbb{P}^n_k,\mathscr{G})^{\vee}.
\end{equation*}
（函子性的）配对 \eqref{eq:hom-trace-pairing} 给出从 \eqref{eq:hom-left-exact} 到 \eqref{eq:cohom-dual-exact} 的一个映射：
\begin{equation*}\tag{29.1.10.4}\label{eq:five-lemma-map}
\begin{tikzcd}[row sep=3.2em, column sep=2.6em]
0 \ar[r] & 0 \ar[r] & \mathrm{H}^n(\mathbb{P}^n_k,\mathscr{F})^{\vee} \ar[r] & \mathrm{H}^n(\mathbb{P}^n_k,\mathscr{E})^{\vee} \ar[r] & \mathrm{H}^n(\mathbb{P}^n_k,\mathscr{G})^{\vee} \\
0 \ar[r] \ar[u, "\alpha"] & 0 \ar[r] \ar[u, "\beta"] & \mathrm{Hom}(\mathscr{F},\omega_{\mathbb{P}^n_k}) \ar[r] \ar[u, "\gamma"] & \mathrm{Hom}(\mathscr{E},\omega_{\mathbb{P}^n_k}) \ar[r] \ar[u, "\delta"] & \mathrm{Hom}(\mathscr{G},\omega_{\mathbb{P}^n_k}). \ar[u, "\epsilon"]
\end{tikzcd}
\end{equation*}
映射 $\alpha$ 与 $\beta$ 显然是\trIsomorphism，而\mbox{练习~29.1.B} 表明 $\delta$ 是一个\trIsomorphism。于是，由五引理的微妙版本（\mbox{练习~1.6.D}；注意到 $\beta$ 与 $\delta$ 均\trInjectiveMap{}而 $\alpha$ \trSurjectiveMap），$\gamma$ \trInjectiveMap。这表明：对\emph{一切}\trCoherentSheaf{} $\mathscr{F}'$，自然映射 $\mathrm{Hom}(\mathscr{F}',\omega_{\mathbb{P}^n_k})\to\mathrm{H}^n(\mathbb{P}^n_k,\mathscr{F}')^{\vee}$ 都\trInjectiveMap；特别地，取 $\mathscr{F}'=\mathscr{G}$ 时也如此。于是 $\epsilon$ \trInjectiveMap。再由五引理微妙版本的对偶命题（由于 $\beta$ 与 $\delta$ 均\trSurjectiveMap，而 $\epsilon$ \trInjectiveMap），$\gamma$ \trSurjectiveMap。\hfill$\square$

\textbf{29.1.11.}\ \emph{数学谜题。}\quad 这里有一个谜题，逼你直面一个可能令人困惑的点。我们将会看到，\mbox{\trDesideratum{}~29.1.1} 对 $\mathbb{P}^1$ 成立，于是 $\omega_{\mathbb{P}^1}\cong\Omega_{\mathbb{P}^1}$。那么\trTraceMap $t\colon\mathrm{H}^1(\mathbb{P}^1,\Omega_{\mathbb{P}^1})\to k$ 究竟是什么？$\mathrm{H}^1(\mathbb{P}^1,\Omega_{\mathbb{P}^1})$ 的\trCech{}复形（取 $\mathbb{P}^1$ 的惯用\trCover）由下式给出
\begin{equation*}\tag{29.1.11.1}\label{eq:cech-complex-omega-p1}
0\longrightarrow\Omega_{\mathbb{P}^1}(U_0)\times\Omega_{\mathbb{P}^1}(U_1)\overset{\alpha}{\longrightarrow}\Omega_{\mathbb{P}^1}(U_0\cap U_1)\longrightarrow 0.
\end{equation*}
若 $U_0=\Spec k[x]$，且 $U_0\cap U_1=\Spec k[x,1/x]$，则 $U_0\cap U_1$ 上的那些\trDifferential{}都形如 $f(x)\,\mathrm{d}x$，其中 $f(x)$ 是一个\trLaurentPolynomial（例如：$(x^{-3}+x^{-1}+3+17x^4)\,\mathrm{d}x$）。要计算 $\mathrm{H}^1(\mathbb{P}^1,\Omega_{\mathbb{P}^1})$，我们需要找出 $U_0\cap U_1$ 上哪些这样的\trDifferential{}落在 \eqref{eq:cech-complex-omega-p1} 中 $\alpha$ 的\trImage{}里。显然，形如 $x^i\,\mathrm{d}x$（$i\ge 0$）的任何一项都能延拓为 $U_0$ 上的一个\trDifferential，从而落在 $\alpha$ 的\trImage{}中。短短一算便可知道，形如 $x^i\,\mathrm{d}x$（$i<-1$）的任何一项都能延拓为 $U_1$ 上的一个\trDifferential。于是 $\alpha$ 的\trCokernel{}可以描述为由 $x^{-1}\,\mathrm{d}x$ 生成的那个一维 $k$-\trVectorSpace。我们有一个到 $k$ 的显然\trIsomorphism：取 $x^{-1}\,\mathrm{d}x$ 的系数，这可以被解释为「取 $0$ 处的留数」。但还有另一个同等好的选择：取 $\infty$ 处的留数——当然，没有任何理由偏爱 $0$ 胜过 $\infty$（或者偏爱 $U_0$ 胜过 $U_1$）！然而这两个留数并\emph{不}相同——它们之和为 $0$（你可以很快算出来——你也可以凭\trRiemannSurface{}理论中的留数定理相信这一点）。那么：哪一个才是迹？

\textbf{29.1.12. \trDesideratum{}：\trDuality{}的一个更强的版本，涉及 $\mathrm{Ext}$。}

\epigraph{你的公式 $\mathrm{H}^{n-p}(X,\mathscr{F})'=\mathrm{Ext}^p_{\OO}(X,F,\Omega^n)$ 令我兴奋不已，因为我深信：无论在解析情形还是代数情形（后者更重要——对我而言），它都是陈述对偶定理的正确方式。}{——J.-P.~\trSerre，致 A.~\trGrothendieck{}书信，1955 年 12 月 22 日 [GrS, p.~21]}

\trVectorBundle{}与 $\mathrm{Hom}$ 的\trSerreDuality{}有一个共同的扩张。为此，我们需要 $\OO$-\trModule{}的 $\mathrm{Ext}$ 概念。我们将在 \S29.2 中发展所需的内容，但眼下只需注意到：若 $X$ 是一个\trRingedSpace，则 $\OO_X$-\trModule{}的\trCategory{}有足够多的\trInjective{}对象（\mbox{定理~23.4.1}），于是我们可以定义 $\mathrm{Ext}^i(\mathscr{F},\cdot)$ 为 $\mathrm{Hom}(\mathscr{F},\cdot)$ 的\trDerivedFunctor（正如 \S23.2.4 中定义的 $A$-\trModule{}的 $\mathrm{Ext}$-函子那样）。

若我们有函子之间的一个\trIsomorphism
\begin{equation*}\tag{29.1.12.1}\label{eq:ext-functorial-iso}
\boxed{\mathrm{Ext}^i(\cdot,\omega_X)\overset{\sim}{\longrightarrow}\mathrm{H}^{n-i}(X,\cdot)^{\vee},}
\end{equation*}
我们就说 $X$ 满足\textbf{$\mathrm{Ext}$ 的函子版本\trSerreDuality}。（$i=0$ 的情形就是 $\mathrm{Hom}$ 的函子版本\trSerreDuality，或者由\mbox{练习~29.1.A}，即迹版本。）

在\mbox{练习~29.2.K} 中，我们将发现：$\mathrm{Ext}$ 的函子版本\trSerreDuality（迹版本亦然）蕴含\trVectorBundle{}的函子版本\trSerreDuality（迹版本亦然）。于是，要证明\mbox{定理~18.5.1}，只须证明 $\mathrm{Ext}$ 的函子版本\trSerreDuality，再加上\mbox{\trDesideratum{}~29.1.1}（见\mbox{注记~29.4.10}）。当 $X$ 为 \trCM{}时，$\mathrm{Ext}$ 的函子版本\trSerreDuality{}将在\mbox{推论~29.3.14} 中得到证明。

\textbf{29.1.13.}\ \emph{$\mathrm{Ext}$ 的\trSerreDuality{}的迹版本。}\quad 与前几种版本一样，$\mathrm{Ext}$ 的函子版本\trSerreDuality{}的函子性也应当有其来源。存在上积
\[
\mathrm{Ext}^a(\mathscr{F},\mathscr{G})\times\mathrm{Ext}^b(\mathscr{E},\mathscr{F})\longrightarrow\mathrm{Ext}^{a+b}(\mathscr{E},\mathscr{F})
\]
（即「\trYoneda{}上积」）。若复合
\[
\resizebox{0.85\textwidth}{!}{\begin{tikzcd}[row sep=3.4em, column sep=1.6em, ampersand replacement=\&]
\mathrm{Ext}^i_X(\mathscr{F},\omega_X)\times\mathrm{Ext}^{n-i}_X(\mathscr{O}_X,\mathscr{F})=\mathrm{Ext}^i_X(\mathscr{F},\omega_X)\times\mathrm{Ext}^{n-i}_X(\mathscr{O}_X,\mathscr{F}) \ar[dr, "\cup"] \& \\
\& \mathrm{Ext}^n_X(\mathscr{O}_X,\omega_X)=\mathrm{H}^n(X,\omega_X) \ar[r, "t"] \& k
\end{tikzcd}}
\]
对一切 $i$ 都是\trPerfectPairing，我们就说 $X$ 满足\textbf{$\mathrm{Ext}$ 的\trSerreDuality{}的迹版本}。我们将无力证明迹版本，因为我们并不会定义这个上积。但一切纯维的\trProper{} \trCM{} $k$-\trScheme{}都满足这一强形式的\trDuality。[Stacks, tag~0FVU]

\textbf{29.1.14. \trCM{}假设的必要性。}\quad 我们指出：凡是在陈述了 \trCM{}假设的地方，这些假设都是必要的。下面的例子在一切情形都适用。设 $X$ 是 $\mathbb{P}^4_k$ 中交于一点 $p$ 的两个二维平面的并。若 $X$ 上存在一个凝聚的\trCanonical{}层 $\omega_X$，则当 $d\gg 0$ 时，由\trSerreDuality{}将会得到 $\mathrm{h}^1(X,\OO_X(-d))=\mathrm{h}^1(X,\omega_X(d))$，而由\trSerre{}消失（\mbox{定理~18.1.3}(ii)）它必须为 $0$。我们将证明事实并非如此。

\textbf{29.1.C.（练习）}\quad 设 $H$ 是 $\mathbb{P}^4$ 中不经过 $p$ 的一个\trHyperplane；比如说，$H$ 就是超平面 $x_0=0$。设 $Z$ 是 $dH$（除子 $x_0^d=0$）与 $X$ 的交。\footnote{译注：原书此处作「for \(d\ge 0\)」；然 $d=0$ 时 $Z=X$ 而 $\mathrm{h}^0(Z,\OO_Z)=1$，疑当为「d > 0」，兹照录原书。——译者注}（用廉价的方式）证明：当 $d\ge 0$ 时，$\mathrm{h}^0(Z,\OO_Z)>1$。利用\trExactSequence
\[
0\longrightarrow\OO_X(-Z)\longrightarrow\OO_X\longrightarrow\OO_Z\longrightarrow 0
\]
来证明 $\mathrm{h}^1(X,\OO_X(-d))>0$。


\section{$\OO$ 模的 $\mathrm{Ext}$ 群与 $\mathrm{Ext}$ 层}

回忆一下：对任意\trRingedSpace{} $X$，范畴 $\mathit{Mod}_{\OO_X}$ 有足够多内射对象（\mbox{定理~23.4.1}）。于是，对 $X$ 上任意 $\OO_X$-\trModule{} $\mathscr{F}$，我们可以定义
\[
\mathrm{Ext}^i_X(\mathscr{F},\cdot)\colon\mathit{Mod}_{\OO_X}\longrightarrow\mathit{Mod}_{\Gamma(X,\OO_X)}
\]
为 $\Hom_X(\mathscr{F},\cdot)$ 的第 $i$ 个右导出函子，并且 $\mathrm{Ext}^i_X(\mathscr{F},\cdot)$ 有相应的\trLongExactSequence。我们类似地定义层版本
\[
\mathscr{E}xt^i_X(\mathscr{F},\cdot)\colon\mathit{Mod}_{\OO_X}\longrightarrow\mathit{Mod}_{\OO_X}
\]
为 $\mathscr{H}om_X(\mathscr{F},\cdot)$ 的第 $i$ 个右导出函子。在这两种情形下，当上下文清楚时，下标 $X$ 常被省略，尽管当讨论中涉及多于一个空间时，这样做可能很危险。（我们在 \S23.2.4 中见过 $A$-\trModule{}的 $\mathrm{Ext}$ 函子。）

\emph{警示：}目前尚不清楚 $\mathit{Mod}_{\OO_X}$ 是否有足够多\trProjectiveObj{}对象（事实上并不成立；见 \S23.4.8），所以我们不能把 $\mathrm{Ext}^i$ 定义成对左变元取导出函子。尽管如此，我们将会看到，它的表现就好像它\emph{是}一个导出函子——它「可用零调对象计算」，并且具有\trLongExactSequence（\mbox{注记~29.2.1}）。

\emph{又一警示：}在这个定义下，如果 $\mathscr{F}$ 与 $\mathscr{G}$ 是一个\trScheme{}上的\trQCSheaf，尚不清楚诸 $\mathscr{E}xt^i(\mathscr{F},\mathscr{G})$ 是否\trQC；而且，\mbox{练习~14.3.A(a)} 中的旁注指出，即便对 $i=0$，这也未必成立。但 \mbox{练习~29.2.F} 会让你安心：只要 $\mathscr{F}$ \trCoh{}，而 $X$ \trLocally\trNoetherian，一切安好。

\mbox{练习~23.5.A}（内射 $\OO_X$-\trModule{}限制到\trOpenSubset{} $\mathfrak{U}\subset X$ 上时，在 $\mathfrak{U}$ 上仍内射）有一系列有用的推论。

\textbf{29.2.A.（练习）}\quad 设 $\mathscr{I}$ 是一个内射 $\OO_X$-\trModule。证明：$\mathscr{H}om_{\OO_X}(\cdot,\mathscr{I})$ 是一个\trExact{}\trContravariantFunctor。（一个相关的事实：$\Hom_{\OO_X}(\cdot,\mathscr{I})$ 是\trExact{}的，由内射性的定义，\mbox{练习~23.2.C(a)}。）

\textbf{29.2.B.（练习）}\quad 设 $X$ 是一个\trRingedSpace，$\mathscr{F}$ 与 $\mathscr{G}$ 是 $\OO_X$-\trModule，而 $\mathfrak{U}$ 是一个\trOpenSubset。请描述一个\trCanonical\trIsomorphism
\[
\mathscr{E}xt^i_X(\mathscr{F},\mathscr{G})|_{\mathfrak{U}}\xrightarrow{\ \sim\ }\mathscr{E}xt^i_{\mathfrak{U}}\bigl(\mathscr{F}|_{\mathfrak{U}},\mathscr{G}|_{\mathfrak{U}}\bigr).
\]
提示：限制函子 $\mathit{Mod}_{\OO_X}\to\mathit{Mod}_{\OO_Y}$ 是\trExact{}的（参看\mbox{练习~2.7.D} 与\mbox{练习~2.7.E}）；限制保持 $\OO$-\trModule{}的内射性（\mbox{练习~23.5.A}）；以及 FHHF 定理（\mbox{练习~1.5.I}）。\footnote{译注：原书此处作「$\mathit{Mod}_{\OO_Y}$」，其中 $Y$ 未在本练习中出现，疑当为 $\mathfrak{U}$（\mbox{练习~2.7.D}／\mbox{练习~2.7.E} 中以 $Y$ 泛指开集），兹照录原书。——译者注}

\textbf{29.2.C.（练习）}\quad 设 $X$ 是一个\trRingedSpace，而 $\mathscr{G}$ 是一个 $\OO_X$-\trModule。

(a) 证明：
\[
\mathscr{E}xt^i(\OO_X,\mathscr{G})=
\begin{cases}
\mathscr{G}\text{，}& \text{若 } i=0\text{；}\\
0\text{，}& \text{其余情形。}
\end{cases}
\]

(b) 请描述一个\trCanonical\trIsomorphism{} $\mathrm{Ext}^i_X(\OO_X,\mathscr{G})\xleftarrow{\ \sim\ }\mathrm{H}^i(X,\mathscr{G})$。

\textbf{29.2.D.（练习）}\quad 用\mbox{练习~29.2.B} 与\mbox{练习~29.2.C(a)} 证明：若 $\mathscr{E}$ 是 $X$ 上的一个\trLocallyFreeSheaf，则当 $i>0$ 时 $\mathscr{E}xt^i(\mathscr{E},\mathscr{G})=0$。

在环上的\trModule{}的范畴中，我们喜欢\trProjectiveObj{}对象甚于内射对象，因为自由模易于打理。若\trScheme{}上的\trLocallyFreeSheaf{}总是\trProjectiveObj{}的，那就再好不过了，但遗憾的是这并不成立。因为若 $\OO_X$ 是投射的，那么当 $i>0$ 时 $\mathrm{Ext}^i(\OO_X,\cdot)$（作为函子）将为零；但由\mbox{练习~29.2.C(b)}，它就是函子 $\mathrm{H}^i(X,\cdot)$，而后者常常非零。尽管如此，我们仍然可以用一个局部自由分解来计算 $\mathscr{E}xt$，正如下面这道练习所示。

\textbf{29.2.E.（重要练习）}\quad 设 $X$ 是一个\trRingedSpace，而
\begin{equation*}
\tag{29.2.0.1}\label{eq:ext-lfr}
\cdots\longrightarrow\mathscr{E}_1\longrightarrow\mathscr{E}_0\longrightarrow\mathscr{F}\longrightarrow 0
\end{equation*}
是 $\OO_X$-\trModule{} $\mathscr{F}$ 的一个由\trLocallyFreeSheaf{}构成的分解。（当然，我们最感兴趣的场合是 $X$ 为\trScheme{}，而 $\mathscr{F}$ \trQC{}、甚至\trCoh{}的情形。）令 $\mathscr{E}_{\bullet}$ 表示 \eqref{eq:ext-lfr} 去掉 $\mathscr{F}$ 后所得的截断。请描述一个\trIsomorphism
\[
\mathscr{E}xt^i(\mathscr{F},\mathscr{G})\xleftarrow{\ \sim\ }\mathrm{H}^i\bigl(\mathscr{H}om(\mathscr{E}_{\bullet},\mathscr{G})\bigr).
\]
换言之，$\mathscr{E}xt^{\bullet}(\mathscr{F},\mathscr{G})$ 可以这样计算：取 $\mathscr{F}$ 的一个局部自由分解，加以截断，施加 $\mathscr{H}om(\cdot,\mathscr{G})$，再取同调。\footnote{译注：原书此处作「$\mathscr{E}xt^{\bullet}(\mathscr{F},\mathscr{G}))$」，右括号多印一枚，当为「$\mathscr{E}xt^{\bullet}(\mathscr{F},\mathscr{G})$」之误，径改。——译者注}提示：选取一个内射分解
\[
0\longrightarrow\mathscr{G}\longrightarrow\mathscr{I}_0\longrightarrow\mathscr{I}_1\longrightarrow\cdots
\]
并考虑以
\[
\begin{tikzcd}[row sep=2.2em, column sep=2.4em,font=\small]
\vdots & \vdots & & \\
\mathscr{H}om(\mathscr{E}_0,\mathscr{I}_1) \ar[r] \ar[u] & \mathscr{H}om(\mathscr{E}_1,\mathscr{I}_1) \ar[r] \ar[u] & \cdots \\
\mathscr{H}om(\mathscr{E}_0,\mathscr{I}_0) \ar[r] \ar[u] & \mathscr{H}om(\mathscr{E}_1,\mathscr{I}_0) \ar[r] \ar[u] & \cdots.
\end{tikzcd}
\]
为 $\mathrm{E}_0$ \trPage{}的\trSpectralSequence。

这个结果很重要：要计算 $\mathscr{E}xt$，我们可以用有限秩的局部自由分解来计算。你可以逐仿射地工作（借助\mbox{练习~29.2.B}），而在每一块仿射上，你可以对左变元使用自由分解。作为\mbox{练习~29.2.E} 的另一条推论：

\textbf{29.2.F.（练习）}

(a) 设 $\mathscr{F}$ 是\trLocally\trNoetherian{}\trScheme{} $X$ 上的一个\trCoherentSheaf，$\mathscr{G}$ 是 $X$ 上的一个\trQCSheaf。证明：$\mathscr{E}xt^i_X(\mathscr{F},\mathscr{G})$ 是 $X$ 上的一个\trQCSheaf。提示：由\mbox{练习~29.2.B}，由于 $\mathscr{E}xt$ 可以在\trOpenCover{}上计算，请证明只需对 $X$ \trAffine{}的情形证明该结果。但在这种情形下，$\mathscr{F}$ 有一个自由分解。（你将会发现 $\mathscr{E}xt^i_X(\mathscr{F},\mathscr{G})(\Spec A)=\mathrm{Ext}^i_A\bigl(\mathscr{F}(\Spec A),\mathscr{G}(\Spec A)\bigr)$。）

(b) 进而，若 $\mathscr{G}$ 也\trCoh，证明 $\mathscr{E}xt^i_X(\mathscr{F},\mathscr{G})$ 也是\trCoh{}的。

\textbf{29.2.1.}\ \emph{\trhRemark。}\quad 陈述「$\mathrm{Ext}^i(\mathscr{F},\mathscr{G})$ 与 $\mathscr{E}xt^i(\mathscr{F},\mathscr{G})$ 在第一个变元上的表现如同导出函子」在好几层意义上都成立。例如，我们有一条相应的\trLongExactSequence，见下一道练习。

\textbf{29.2.G.（练习）}\quad 设 $0\to\mathscr{F}'\to\mathscr{F}\to\mathscr{F}''\to 0$ 是\trRingedSpace{} $X$ 上诸 $\OO_X$-\trModule{}的一个\trExactSequence。对任意 $\OO_X$-\trModule{} $\mathscr{G}$，请描述一条\trLongExactSequence
\[
\begin{gathered}
0\longrightarrow\Hom(\mathscr{F}'',\mathscr{G})\longrightarrow\Hom(\mathscr{F},\mathscr{G})\longrightarrow\Hom(\mathscr{F}',\mathscr{G})\\
\longrightarrow\mathrm{Ext}^1(\mathscr{F}'',\mathscr{G})\longrightarrow\mathrm{Ext}^1(\mathscr{F},\mathscr{G})\longrightarrow\mathrm{Ext}^1(\mathscr{F}',\mathscr{G})\longrightarrow\cdots.
\end{gathered}
\]
（于是，$\mathrm{Ext}^i(\cdot,\mathscr{G})$ 构成一个反变\trDeltaFunctor——如 \S23.2.7 中所定义——尽管它并不是作为导出函子来定义的。）请对 $\mathscr{E}xt$ 证明类似的结果。提示：取一个内射分解 $0\to\mathscr{G}\to\mathscr{I}^0\to\cdots$。利用「若 $\mathscr{I}$ 内射，则 $\Hom(\cdot,\mathscr{I})$ \trExact」（内射性的定义，\mbox{练习~23.2.C(a)}）这一事实。由此得到\trComplexN{}的一个\trShortExactSequence
\[
0\longrightarrow\Hom(\mathscr{F}'',\mathscr{I}^{\bullet})\longrightarrow\Hom(\mathscr{F},\mathscr{I}^{\bullet})\longrightarrow\Hom(\mathscr{F}',\mathscr{I}^{\bullet})\longrightarrow 0
\]
并取上同调中的\trLongExactSequence。

\textbf{29.2.2. 对 $\mathrm{Ext}$ 的函子性 \trSerreDuality{}在\trProjectiveSpace{}上成立。}

我们现在证明：对 $\mathrm{Ext}$ 的函子性 \trSerreDuality{}在\trProjectiveSpace{}上成立。我们将使用泛\trDeltaFunctor{}的机制（\S23.2.6），所以你既不妨快速浏览一下那一节，也可以干脆忽略这段讨论。

\textbf{29.2.H.（练习）}\quad 证明：$\bigl(\mathrm{Ext}^i_{\mathbb{P}^n_k}(\cdot,\omega_{\mathbb{P}^n_k})\bigr)$ 是一个（反变的）泛\trDeltaFunctor。提示：$\mathrm{Ext}$ 在其第一个变元上不是导出函子，所以你不能用\mbox{推论~23.2.11} 的「\trProjectiveObj」版本。作为替代，请使用\mbox{定理~23.2.9}，以及「对每个 $\mathscr{F}$ 都存在满态射 $\OO(m)^{\oplus N}\to\mathscr{F}$，其中 $m<0$」这一事实。

\textbf{29.2.I.（练习）}\quad 证明：$\bigl(\mathrm{H}^{n-i}(\mathbb{P}^n_k,\cdot)^{\vee}\bigr)$ 是一个泛\trDeltaFunctor。（诸 $\delta$-映射是什么？）提示：试用与上一道练习相同的想法。

\mbox{命题~29.1.10} 给出函子之间的一个\trIsomorphism{} $\mathrm{Ext}^0_{\mathbb{P}^n_k}(\cdot,\omega_{\mathbb{P}^n_k})\xleftarrow{\ \sim\ }\mathrm{H}^n(\mathbb{P}^n_k,\cdot)^{\vee}$，于是由泛\trDeltaFunctor{}的\mbox{定义~23.2.8}，我们便有诸\trDeltaFunctor{}之间的一个\trIsomorphism{} $\bigl(\mathrm{Ext}^i_{\mathbb{P}^n_k}(\cdot,\omega_{\mathbb{P}^n_k})\bigr)\xleftarrow{\ \sim\ }\bigl(\mathrm{H}^{n-i}(\mathbb{P}^n_k,\cdot)^{\vee}\bigr)$，从而证明了对 $\mathbb{P}^n_k$ 而言对 $\mathrm{Ext}$ 的函子性 \trSerreDuality。

\textbf{29.2.3. 两道有用的练习，以及\trProjectiveSpace{}上对\trVectorBundle{}的函子性 \trSerreDuality。}

\textbf{29.2.J.（练习）}\quad 设 $X$ 是一个\trRingedSpace，$\mathscr{F}$ 与 $\mathscr{G}$ 是 $\OO_X$-\trModule，而 $\mathscr{E}$ 是 $X$ 上的一个\trLocallyFreeSheaf。请描述诸\trIsomorphism
\[
\mathscr{E}xt^i(\mathscr{F}\otimes\mathscr{E}^{\vee},\mathscr{G})\xleftarrow{\ \sim\ }\mathscr{E}xt^i(\mathscr{F},\mathscr{G}\otimes\mathscr{E})\xleftarrow{\ \sim\ }\mathscr{E}xt^i(\mathscr{F},\mathscr{G})\otimes\mathscr{E}
\]
以及
\[
\mathrm{Ext}^i(\mathscr{F}\otimes\mathscr{E}^{\vee},\mathscr{G})\xleftarrow{\ \sim\ }\mathrm{Ext}^i(\mathscr{F},\mathscr{G}\otimes\mathscr{E}).
\]
提示：证明若 $\mathscr{I}$ 内射，则 $\mathscr{I}\otimes\mathscr{E}$ 也内射。

\textbf{29.2.K.（练习）}\quad 若 $\mathscr{F}$ 是\trScheme{} $X$ 上的一个\trLocallyFreeSheaf，而 $\mathscr{G}$ 是 $X$ 上任意一个\trCoherentSheaf，请描述一个\trIsomorphism{} $\mathrm{Ext}^i(\mathscr{F},\mathscr{G})\xleftarrow{\ \sim\ }\mathrm{H}^i(X,\mathscr{F}^{\vee}\otimes\mathscr{G})$（对 $\mathscr{F}$ 函子）。取 $\mathscr{G}$ 为 $\omega_X$，证明：对 $\mathrm{Ext}$ 的函子性 \trSerreDuality{}蕴含对\trVectorBundle{}的函子性 \trSerreDuality。（你也不妨琢磨一下迹版本。）于是，由 \S29.2.2，对\trVectorBundle{}的函子性 \trSerreDuality{}在\trProjectiveSpace{}上成立。提示：\mbox{练习~29.2.J} 与\mbox{练习~29.2.C(b)}。

\textbf{29.2.4. $\mathrm{Ext}$ 的局部到整体\trSpectralSequence。}

「层」$\mathscr{E}xt$ 与「整体」$\mathrm{Ext}$ 之间由一个\trSpectralSequence{}相联系。一旦我们证明了 $\mathscr{H}om(\mathscr{F},\mathscr{I})$ 对函子 $\Gamma$ 是\trAcyclic{}的，这就是\trGrothendieck{}函子复合谱序列（\mbox{定理~23.3.5}）的直接应用。

\textbf{29.2.L.（练习）}\quad 设 $\mathscr{I}$ 是一个内射 $\OO_X$-\trModule。证明：$\mathscr{H}om(\mathscr{F},\mathscr{I})$ 是\trFlasque{}的（\mbox{定义~23.4.6}），从而由\mbox{练习~23.4.I} 是 $\Gamma$-\trAcyclic{}的。提示：设 $j\colon\mathfrak{U}\hookrightarrow V$ 是\trOpenSubset{}之间的一个包含。我们希望证明 $\mathscr{H}om(\mathscr{F},\mathscr{I})(V)\to\mathscr{H}om(\mathscr{F},\mathscr{I})(\mathfrak{U})$ 是\trSurjectiveMap。注意 $\mathscr{I}|_V$ 在 $V$ 上内射（\mbox{练习~23.5.A}）。把\trExact{}函子 $\Hom_V(\cdot,\mathscr{I}|_V)$ 施加到 $V$ 上诸\trSheaf{}之间的包含态射 $j_!(\mathscr{F}|_{\mathfrak{U}})\hookrightarrow\mathscr{I}|_V$ 上（\S23.4.7）。

\textbf{29.2.M.（练习：$\mathrm{Ext}$ 的局部到整体\trSpectralSequence）}\quad 设 $X$ 是一个\trRingedSpace，而 $\mathscr{F}$ 与 $\mathscr{G}$ 是 $\OO_X$-\trModule。请描述一个以 $\mathrm{H}^i\bigl(X,\mathscr{E}xt^j(\mathscr{F},\mathscr{G})\bigr)$ 为 $\mathrm{E}_2$ 项、\trAbut{}于 $\mathrm{Ext}^{i+j}(\mathscr{F},\mathscr{G})$ 的\trSpectralSequence。（提示：使用\trGrothendieck{}函子复合谱序列，\mbox{定理~23.3.5}。回忆 $\Hom(\mathscr{F},\cdot)=\Gamma\bigl(\mathscr{H}om(\mathscr{F},\cdot)\bigr)$，\mbox{练习~2.3.C}。）


\section{\trProjective{} $k$-\trScheme{}的\trSerreDuality}

凭借\trProjectiveSpace{}的\trSerreDuality，我们现在来证明\trProjective{} $k$-\trScheme{}的\trSerreDuality{}的种种版本。

\textbf{29.3.1. 回顾古怪的函子 $\pi^{!?}$。}

关键的构造是\trPushforward{} $\pi_*$ 的一个\emph{\trRightAdjoint} $\pi^{!?}$（当 $\pi$ 是一个\trFinite\trMorphism{}时），见 \S17.1.6 中的讨论。若你尚未读过那一节，现在应当读一读。

\textbf{29.3.2.}\ \emph{警示。}\quad 我们只对 $Y$ 上的\trQCSheaf{}定义过 $\pi^{!?}$，并非对一切 $\OO_Y$-\trModule。若 $\mathscr{G}$ 是一个 $\OO_X$-\trModule{}（不必\trQC），那么一般说来，如何为这个构造赋予意义、从而定义出一个 $\OO_Y$-\trModule{} $\pi^{!?}\mathscr{G}$，并\emph{不}清楚。（试试看！）然而，在 $\pi$ 是\trClosedEmbedding{}的特殊情形下，我们\emph{能够}理解 $\pi^{!?}\mathscr{G}$，如\mbox{练习~29.4.A} 所描述的那样。

\textbf{29.3.3.}\ \emph{注记。}\quad 若 $\pi$ \trFinite{}且\trFlat——这正是我们最感兴趣的情形——则 $\pi^{!?}$ 与 $\pi^!$ 相合（就「相合」一词唯一可能的含义而言），而后者是\trGrothendieck{}「六运算」（six operations）之一；我们的记号也正源于此（「$\pi^!$ 的一个奇怪变体」）。但 $\pi^!$ 天然地生活在\trDerivedCategory{}的世界里，所以我们在此不讨论它。（关于 $\pi^!$ 的更多内容，见 [KS1, Ch.~III]。）

现在我们把 $\pi^{!?}$ 的这套机制应用于\trSerreDuality。

\textbf{29.3.4. 射影 $n$ 维概形是 $\mathbb{P}^n$ 的\trFinite\trCover（finite cover）。}

\textbf{29.3.5. \trhProposition。}\ ——\ \emph{设 $X$ 是\trDimension{}至多为 $n$ 的\trProjective{} $k$-\trScheme。}

\emph{(a)~存在一个\trFinite\trMorphism{} $\pi\colon X\to\mathbb{P}^n$。}

\emph{(b)~若进而 $X$ 是\trPureDimension{}为 $n$ 的 \trCM，则 $\pi$ 是\trFlat{}的。}

\textbf{29.3.6.}\ \emph{致专家的注记。}\quad 我们在这里刻意允许 $n$ 大于 $\dim X$，这不是因为它有用，而是要指出：这个证明是行得通的，而一旦假设不当，我们证得的便是某种奇怪的东西。当 $\dim X=n$ 而 $X$ 并非 \trCM{}时，\mbox{命题~29.3.5(a)} 仍然有趣。而当 $\dim X<n$ 时，\mbox{命题~29.3.5(a)} 又蕴涵什么呢？

\emph{证明。}\quad (b) 部分由 (a) 部分与\trMiracleFlat{}定理~26.2.11 得出，所以我们来证 (a)。

选取一个\trClosedEmbedding{} $j\colon X\hookrightarrow\mathbb{P}^N$。为叙述简明，先设 $k$ 是无限域。由\mbox{练习~12.3.D(d)}，存在一个余维数为 $n+1$（比补维数恰小 $1$）且与 $X$ 不相交的\trLinearSpace{} $L$。从 $L$ 作\trProjection{}便给出一个态射 $\pi\colon X\to\mathbb{P}^n$。态射 $\pi$ 是\trAffine{}的（$\mathbb{P}^N\setminus L\to\mathbb{P}^n$ 是\trAffine{}的，而\trClosedEmbedding{} $X\hookrightarrow\mathbb{P}^N\setminus L$ 同一切\trClosedEmbedding{}一样是\trAffine{}的），又是\trProjective{}的（由射影态射的消去定理，\mbox{练习~17.3.C}；或者更直接地由\mbox{练习~17.3.D}），于是 $\pi$ 是\trFinite{}的（局部\trNoether{}概形上的射影\trAffine{}态射是\trFinite{}的，\mbox{推论~18.1.5}）。

\textbf{29.3.A.（练习）}\quad 在不假定 $k$ 无限的情形下证明\mbox{命题~29.3.5(a)}。提示：可先用\mbox{练习~12.3.D(c)} 证明：存在某个 $d$，使得有 $n+1$ 张 $d$ 次\trHypersurface{}的交与 $X$ 不相交。然后把上述论证应用于 $\mathbb{P}^N$ 的 $d$ 次\trVeronese{}嵌入（\S9.3.6）即可。\hfill$\square$

\textbf{29.3.7.}\ \emph{注记。}\quad 更一般地，上述论证表明：任何维数为 $n$（不必纯维）的\trProjective{} $k$-\trScheme{}都能表示为 $\mathbb{P}^n$ 的一个\trFinite\trCover。这或许可以看作\trNoetherNormalization{}的一个\trProjective{}版本。

\textbf{29.3.8. 借助 $\pi^{!?}$ 的 $X$ 上的\trSerreDuality。}

设 $\pi\colon X\to Y$ 是\trPureDimension{}为 $n$ 的\trProjective{} $k$-\trScheme{}之间的一个\trFinite\trMorphism。（我们马上就会在 $Y=\mathbb{P}^n_k$ 的情形下应用它，但我们不妨避开干扰，也无须那样特化。）

\textbf{29.3.9. \trhProposition。}\ ——\ \emph{设 $\pi$ \trFlat。若\trVectorBundle{}的函子版本\trSerreDuality{}对 $Y$ 成立，带\trCanonical{}层 $\omega_Y$，则\trVectorBundle{}的函子版本\trSerreDuality{}对 $X$ 成立，带\trCanonical{}层 $\pi^{!?}(\omega_Y)$。}

\emph{证明。}\quad 对每个 $i$ 以及 $X$ 上每个有限秩的\trLocallyFreeSheaf{} $\mathscr{F}$，我们有诸\trIsomorphism{}（随 $\mathscr{F}$ 函子）：
\[
{\small
\begin{aligned}
\mathrm{H}^{n-i}(X,\mathscr{F}^{\vee}\otimes\pi^{!?}\omega_Y)\cong{}& \mathrm{H}^{n-i}\bigl(Y,\pi_*(\mathscr{F}^{\vee}\otimes\pi^{!?}\omega_Y)\bigr) && \text{（$\pi$ 的\trAffineness，\S18.1~(v)）}\\
\cong{}& \mathrm{H}^{n-i}\bigl(Y,\pi_*(\mathscr{H}om_X(\mathscr{F},\pi^{!?}\omega_Y))\bigr) && \text{（\mbox{练习~14.1.G}）}\\
\cong{}& \mathrm{H}^{n-i}\bigl(Y,\mathscr{H}om_Y(\pi_*\mathscr{F},\omega_Y)\bigr) && \text{（\mbox{式~(17.1.8.2)}）}\\
\cong{}& \mathrm{H}^{n-i}\bigl(Y,(\pi_*\mathscr{F})^{\vee}\otimes\omega_Y\bigr) && \text{（\mbox{练习~24.3.E} 与\mbox{练习~14.1.G}）}\\
\cong{}& \mathrm{H}^{i}(Y,\pi_*\mathscr{F})^{\vee} && \text{（$Y$ 上\trVectorBundle{}的\trSerreDuality）}\\
\cong{}& \mathrm{H}^{i}(X,\mathscr{F})^{\vee}. && \text{（$\pi$ 的\trAffineness，\S18.1~(v)）}
\end{aligned}}
\]
（\trFlatness{}假设用于证明 $\pi_*\mathscr{F}$ 是\trLocallyFree{}的。）\hfill$\square$

\textbf{29.3.10. \trhCorollary。}\ ——\ \emph{\trVectorBundle{}的函子版本\trSerreDuality{}对一切纯维的 \trCM{} \trProjective{} $k$-\trScheme{}成立。}

\emph{证明。}\quad 把\mbox{命题~29.3.9} 与\mbox{命题~29.3.5(b)} 以及 \S29.2.2 结合起来。\hfill$\square$

\textbf{29.3.11. \trhProposition。}\ ——\ \emph{设 $\pi\colon X\to Y$ 是\trProjective{} $k$-\trScheme{}之间的一个\trFinite\trMorphism。若 $\mathrm{Hom}$ 的\trSerreDuality{}对 $Y$ 成立，带\trCanonical{}层 $(\omega_Y,t_Y)$，则 $\mathrm{Hom}$ 的\trSerreDuality{}对 $X$ 成立，带\trCanonical{}层 $\pi^{!?}\omega_Y$。}

回忆一下\mbox{练习~29.1.A}：对 $\mathrm{Hom}$ 的\trSerreDuality{}而言，迹版本与函子版本等价。因此我们径直使用「$\mathrm{Hom}$ 的\trSerreDuality」这个说法。

注意，我们对 $\pi$ 不作任何\trFlat{}假设（不同于\trVectorBundle{}的情形，也不同于稍后 $\mathrm{Ext}$ 的情形中那些相应的命题）。

\emph{证明。}\quad 我们有函子性的\trIsomorphism
\[
\begin{aligned}
\Hom_X(\mathscr{F},\pi^{!?}\omega_Y)\cong{}& \Hom_Y(\pi_*\mathscr{F},\omega_Y) && \text{（$(\pi_*,\pi^{!?})$ 的\trAdjointness）}\\
\cong{}& \mathrm{H}^n(Y,\pi_*\mathscr{F})^{\vee} && \text{（$Y$ 的 $\mathrm{Hom}$ \trSerreDuality）}\\
\cong{}& \mathrm{H}^n(X,\mathscr{F})^{\vee}. && \text{（$\pi$ 的\trAffineness）}
\end{aligned}
\]
\hfill$\square$

你可以把上述论证展开，用 $Y$ 的\trTraceMap{} $t_Y$ 表示出 $X$ 的\trTraceMap{} $t_X$：它就是复合
\[
\mathrm{H}^n(X,\pi^{!?}\omega_Y)\longrightarrow\mathrm{H}^n(Y,\pi_*\pi^{!?}\omega_Y)\longrightarrow\mathrm{H}^n(Y,\omega_Y)\overset{t_Y}{\longrightarrow} k.
\]

\textbf{29.3.12. \trhCorollary。}\ ——\ \emph{$\mathrm{Hom}$ 的\trSerreDuality{}的函子版本与迹版本对一切纯维\trProjective{} $k$-\trScheme{}成立（不带任何 \trCM{}假设）。}

\emph{证明。}\quad 把\mbox{命题~29.3.11} 与\mbox{命题~29.3.5(a)} 以及\mbox{命题~29.1.10} 结合起来。\hfill$\square$

\textbf{29.3.13. \trhProposition。}\ ——\ \emph{若 $X$ 是\trPureDimension{}为 $n$ 的 \trCM{} \trProjective{} $k$-\trScheme，而 $\pi\colon X\to\mathbb{P}^n$ 是一个\trFlat{}\trFinite\trMorphism（如\mbox{命题~29.3.5(b)} 中那样），则 $\mathrm{Ext}$ 的函子版本\trSerreDuality{}对 $X$ 成立，带\trCanonical{}层 $\omega_X:=\pi^{!?}\omega_{\mathbb{P}^n}$。}

与\mbox{命题~29.3.9} 和\mbox{命题~29.3.11} 不同的是，我们并没有对一般的 $\pi\colon X\to Y$ 陈述这条命题；我们将用到「目标是 $\mathbb{P}^n$」这一事实。

\emph{证明。}\quad 论证将与 \S29.2.2 平行：我们将给出诸\trDeltaFunctor{}之间的一个\trIsomorphism
\begin{equation*}\tag{29.3.13.1}\label{eq:ext-serre-duality-iso}
\mathrm{Ext}^i_X(\mathscr{F},\pi^{!?}\omega_{\mathbb{P}^n})\overset{\sim}{\longrightarrow}\mathrm{H}^{n-i}(X,\mathscr{F})^{\vee}.
\end{equation*}
$i=0$ 时的\trIsomorphism{}我们已经有了（\mbox{推论~29.3.12}），所以只须证明 \eqref{eq:ext-serre-duality-iso} 的两侧都是泛\trDeltaFunctor（\S23.2.6）。为此，我们来验证\mbox{定理~23.2.9} 的判据。对 $X$ 上任意\trCoherentSheaf{} $\mathscr{F}$，我们可以找到满态射 $\OO(-m)^{\oplus N}\to\mathscr{F}$（其中 $m\gg 0$，而 $N$ 为某个数），于是只须证明：当 $m\gg 0$ 时，
\[
\mathrm{Ext}^i_X(\OO(-m),\pi^{!?}\omega_{\mathbb{P}^n})=0\qquad\text{与}\qquad \mathrm{H}^{n-i}(X,\OO(-m))^{\vee}=0
\]
对 $i>0$ 与 $m\gg 0$ 成立。对于前者，我们有（当 $i>0$ 且 $m\gg 0$ 时）
\[
\begin{aligned}
\mathrm{Ext}^i_X(\OO(-m),\pi^{!?}\omega_{\mathbb{P}^n})&\cong\mathrm{H}^i\bigl(X,(\pi^{!?}\omega_{\mathbb{P}^n})(m)\bigr) && \text{（\mbox{练习~29.2.K}）}\\
&=0 && \text{（$X$ 的\trSerre{}消失，\mbox{定理~18.1.3(ii)}）。}
\end{aligned}
\]
对于后者，我们有
\[
\begin{aligned}
\mathrm{H}^{n-i}(X,\OO_X(-m))&=\mathrm{H}^{n-i}(X,\OO_X\otimes\pi^*\OO(-m))\\
&\cong\mathrm{H}^{n-i}\bigl(\mathbb{P}^n,(\pi_*\OO_X)\otimes\OO(-m)\bigr)
\end{aligned}
\]
（由\trProjection{}公式，\mbox{练习~18.7.E(b)}）。而 $\pi_*\OO_X$ 是\trLocallyFree{}的（\mbox{练习~24.3.E}），于是由 $\mathbb{P}^n$ 上\trVectorBundle{}的函子版本\trSerreDuality，我们有（当 $m\gg 0$ 时）
\[
\begin{aligned}
\mathrm{H}^{n-i}\bigl(\mathbb{P}^n,(\pi_*\OO_X)\otimes\OO(-m)\bigr)\cong{}& \mathrm{H}^{i}\bigl(\mathbb{P}^n,(\pi_*\OO_X)^{\vee}\otimes\OO(m)\otimes\omega_{\mathbb{P}^n}\bigr)^{\vee}\\
\cong{}& \mathrm{H}^{i}\bigl(\mathbb{P}^n,((\pi_*\OO_X)^{\vee}\otimes\omega_{\mathbb{P}^n})\otimes\OO(m)\bigr)^{\vee}\\
={}& 0 \qquad\text{（$\mathbb{P}^n$ 的\trSerre{}消失，\mbox{定理~18.1.3(ii)}）。}
\end{aligned}
\]
\hfill$\square$

\textbf{29.3.14. \trhCorollary。}\ ——\ \emph{$\mathrm{Ext}$ 的函子版本\trSerreDuality{}对一切纯维的 \trCM{} \trProjective{} $k$-\trScheme{}成立。}

\emph{证明。}\quad 把\mbox{命题~29.3.13} 与\mbox{命题~29.3.5(b)} 以及 \S29.2.2 的射影空间情形结合起来。\hfill$\square$

\textbf{29.3.15.}\ \extrahard\ \emph{\trhRemark：\trVectorBundle{}的\trSerreDuality{}的迹版本。}\quad 下面是一个梗概，说明如何把上述证明加以扩展，以对一切纯维的 \trCM{} \trProjective{} $k$-\trScheme{}证明\trVectorBundle{}的\trSerreDuality{}的迹版本。沿用\mbox{命题~29.3.13} 的记号，\trYoneda{}上积（我们尚未定义它）继以\trTraceMap，便给出一个配对
\[
\mathrm{Ext}^i_X(\mathscr{F},\omega_X)\times\mathrm{H}^{n-i}(X,\mathscr{F})\overset{\cup}{\longrightarrow}\mathrm{H}^n(X,\omega_X)\overset{t}{\longrightarrow} k\,,
\]
它随 $\mathscr{F}$ 函子。这诱导出 \eqref{eq:ext-serre-duality-iso} 的一个函子性的映射。可以证明，这个特定的映射是一个\trDeltaFunctor{}态射。于是由上述论证，这个特定的映射是一个\trIsomorphism。

\textbf{29.3.16.}\quad 正如 \S29.1.13 中所描述的，实际情形还要更好：纯维的\trProper{} $k$-\trScheme{}满足最好的那种\trSerreDuality，即 $\mathrm{Ext}$ 的迹版本。[Stacks, tag~0FVU]

\textbf{29.3.17. 结论。}\quad 本节中的这些证明或许出人意料地容易。$\pi^{!?}$ 一经定义，我们就只用了「它是 $\pi_*$ 的右伴随」这唯一一桩事实。我们付出的代价是：从这个证明里，我们对\trCanonical{}层究竟是什么模样几乎一无所知。例如，若 $X$ \trSmooth，并不能明白 $\omega_X$ 是不是一个\trLineBundle！（我们将在下一节中补救这一点。）

你或许也已经感觉到，这些想法如何通向一个更大的故事：为了在更一般的情形下定义并证明\trDuality，我们正试图找出一个函子 $\pi^!$，使之在我们力所能及的最大普遍性上成为 $\pi_*$ 的\trRightAdjoint。这将引导我们在比\trCoherentSheaf{}更一般的框架中工作——\trCoherentSheaf{}的\trDerivedCategory{}（乃至更远）。


\section{$\omega_X$ 的\trAdjunctionFormula{}与 $\omega_X=\mathscr{K}_X$}

（在本节——全书最后一节——中，尽管我们曾在 \S29.1.6 中就术语提出过异议，我们仍将把「\trDualizingSheaf」一词用于 \trSerreDuality{}的\trCanonical{}层，原因在于：本节的全部要点正是要表明，「\trCanonical{}层」这一短语的两种用法在\trSmooth\trProjective\trVariety{} $X$ 的情形下指的是同一个东西：在 \S29.1.5 的 $\mathrm{Hom}$ 的\trSerreDuality{}中出现的那个层 $\omega_X$，以及 $\mathscr{K}_X=\det\Omega_X$。）

\trDualizingSheaf{}在用\trEffectiveCartierDivisor{}作\trSlicing{}时表现良好，而这对得到\mbox{\trDesideratum{}~29.1.1}——亦即 $\mathscr{K}=\det\Omega$ 是\trSerre{}对偶化的这一奇迹（\S18.5.2）——将会很有用。为了说明这一点，我们将对满足 $\mathrm{Ext}$ 的\trSerreDuality{}的\trVariety{}之\trSubvariety，给出其\trDualizingSheaf{}的一个描述。

作为热身，我们将在 $\pi$ 为\trClosedEmbedding{}的特殊情形下，把 $\pi^!$ 的概念推广到 $\OO$-\trModule{}上。

\textbf{29.4.A.（练习）}\quad 设 $\pi\colon X\to Y$ 是概形之间的一个\trClosedEmbedding，而 $\mathscr{G}$ 是一个 $\OO_Y$-\trModule。请解释为什么 $\mathscr{H}om_Y(\pi_*\OO_X,\mathscr{G})$ 自然地带有 $\OO_X$-\trModule{}的结构。提示：若 $\mathscr{I}$ 是 $X$ 的\trIdealSheaf，请解释 $\mathscr{H}om_Y(\pi_*\OO_X,\mathscr{G})$（在某个\trOpenSubset{} $\mathfrak{U}\subset Y$ 上）如何被「在 $X$ 上消失的函数」（$\mathscr{I}(\mathfrak{U})$ 中的元素）所零化。

\textbf{29.4.1.}\ \emph{\trhDefinition{}。}\quad 于是，若 $\pi\colon X\to Y$ 是一个\trClosedEmbedding，我们就已经定义了一个映射 $\pi^!\colon\mathit{Mod}_{\OO_Y}\to\mathit{Mod}_{\OO_X}$，它是\mbox{练习~17.1.K(a)} 中那个映射的扩张。在这种 $\pi$ 为\trClosedEmbedding{}的情形下，记 $\pi^!\mathscr{G}:=\pi^{-1}\mathscr{H}om(\pi_*\OO_X,\mathscr{G})$ 是合理且正确的，也就是说，在定义中使用符号 $\pi^{-1}$。

\textbf{29.4.B.（练习）}\quad 我们仍处于\mbox{练习~29.4.A} 的情形中。

(a) 证明：$(\pi_*,\pi^!)$ 是 $\mathit{Mod}_{\OO_X}$ 与 $\mathit{Mod}_{\OO_Y}$ 之间的一个\trAdjointPair。

(b) 证明：$\pi^!$ 把内射的 $\OO_Y$-\trModule{}送到内射的 $\OO_X$-\trModule。（提示：$\pi_*$ 是\trExact{}的；用\mbox{练习~23.5.B}。）

我们现在把这一讨论从 $\mathscr{H}om$ 推广到 $\mathscr{E}xt$。设 $\pi\colon X\to Y$ 是一个\trClosedEmbedding。那么，对任意 $\OO_Y$-\trModule{} $\mathscr{F}$，$\mathscr{E}xt^i_Y(\pi_*\OO_X,\mathscr{F})$ 自然地带有 $\OO_X$-\trModule{}的结构。（理由：我们通过如下方式计算它：取 $\mathscr{F}$ 的一个由 $\OO_Y$-\trModule{}构成的内射分解，加以截断，然后施加 $\mathscr{H}om_Y(\pi_*\OO_X,\cdot)$。而对任意 $\OO_Y$-\trModule{} $\mathscr{G}$，$\mathscr{H}om_Y(\pi_*\OO_X,\mathscr{G})$ 就带有 $\OO_X$-\trModule{}的结构，也就是 $\pi^!\mathscr{G}$。）为了强调它作为 $\OO_X$-\trModule{}的结构，我们把它写作 $\pi^{-1}\mathscr{E}xt^i_Y(\pi_*\OO_X,\mathscr{F})$。

\textbf{29.4.C.（练习）}\quad 证明：若 $Y$（从而 $X$）是\trLocally\trNoetherian{}的，而 $\mathscr{G}$ 是 $Y$ 上的一个\trCoherentSheaf，则 $\pi^{-1}\mathscr{E}xt^i_Y(\pi_*\OO_X,\mathscr{G})$ 是 $X$ 上的一个\trCoherentSheaf。（提示：$\mathscr{E}xt^i_Y(\pi_*\OO_X,\mathscr{G})$ 是 $Y$ 上的\trCoherentSheaf，\mbox{练习~29.2.F(b)}。）

\textbf{29.4.2. 用环绕簇的对偶化层表示子簇的对偶化层。}\quad 在本节余下的部分中，$\pi\colon X\to Y$ 都将是纯维\trProjective{} $k$-\trScheme{}之间的一个\trClosedEmbedding，二者的\trDimension{}分别为 $n$ 与 $N$，其中 $Y$ 满足对 $\mathrm{Ext}$ 的函子性\trSerreDuality。令 $r=N-n$（即 $X$ 在 $Y$ 中的\trCodimension）。由\mbox{推论~29.3.12}，$X$ 满足对 $\mathrm{Hom}$ 的函子性\trSerreDuality。我们现在要用 $\omega_Y$ 来识别 $\omega_X$。

\textbf{29.4.3. \trhTheorem。}\ ——\ \emph{我们有 $\omega_X=\pi^{-1}\mathscr{E}xt^r_Y(\pi_*\OO_X,\omega_Y)$。}

在证明\mbox{定理~29.4.3} 之前，我们先来解释，那些「先前的」$\mathscr{E}xt^i_Y(\pi_*\OO_X,\omega_Y)$ 都发生了什么。

\textbf{29.4.4. \trhProposition。}\ ——\ \emph{设 $\pi\colon X\hookrightarrow Y$ 是纯维\trProjective{} $k$-\trScheme{}之间的一个\trClosedEmbedding，二者的\trDimension{}分别为 $n$ 与 $N$，且 $Y$ 满足对 $\mathrm{Ext}$ 的函子性\trSerreDuality。则对一切 $i<r:=N-n$，都有 $\mathscr{E}xt^i_Y(\pi_*\OO_X,\omega_Y)=0$。}

\emph{证明。}\quad 由于 $\mathscr{E}xt^i_Y(\pi_*\OO_X,\omega_Y)$ 是\trCoh{}的（\mbox{练习~29.2.F}），故只须证明：当 $m\gg 0$ 时，$\mathscr{E}xt^i_Y(\pi_*\OO_X,\omega_Y)\otimes\OO(m)$ 没有非零的整体截面（因为对 $Y$ 上任意\trCoherentSheaf{} $\mathscr{G}$，当 $m\gg 0$ 时 $\mathscr{G}(m)$ 由整体截面生成——见\trSerre{}定理~A，\mbox{定理~16.1.5}）。由\mbox{练习~29.2.J}，
\[
\mathrm{H}^j(Y,\mathscr{E}xt^i_Y(\pi_*\OO_X,\omega_Y)(m))=\mathrm{H}^j(Y,\mathscr{E}xt^i_Y(\pi_*\OO_X,\omega_Y(m))).
\]
若 $j>0$，则当 $m\gg 0$ 时，由\trSerre{}消失，$\mathrm{H}^j(Y,\mathscr{E}xt^i_Y(\pi_*\OO_X,\omega_Y)(m))=0$。于是，由对 $\mathrm{Ext}$ 的局部到整体\trSpectralSequence（\mbox{练习~29.2.M}），
\[
\mathrm{H}^0(Y,\mathscr{E}xt^i_Y(\pi_*\OO_X,\omega_Y(m)))=\mathrm{Ext}^i_Y(\pi_*\OO_X,\omega_Y(m)).
\]
再用一次\mbox{练习~29.2.J}，继而由 $Y$ 上对 $\mathrm{Ext}$ 的函子性\trSerreDuality，
\[
\mathrm{Ext}^i_Y(\pi_*\OO_X,\omega_Y(m))=\mathrm{Ext}^i_Y(\pi_*\OO_X(-m),\omega_Y)=\mathrm{H}^{N-i}(Y,\pi_*\OO_X(-m)),
\]
它在 $N-i>n$ 时为 $0$，因为\trProjective\trScheme{}上一个\trQCSheaf{}的\trCohomology{}在高于该层\trSupport{}之\trDimension{}的次数上消失（维数\trCohomology{}消失，\mbox{定理~18.2.6}）。\footnote{译注：原书上述第二、三式中的「$\mathrm{Ext}^i_Y(\pi_*\OO_X,\omega_Y(m))$」均多印一枚右括号（作「$\mathrm{Ext}^i_Y(\pi_*\OO_X,\omega_Y(m)))$」），当为排印之误，径改。——译者注}\hfill$\square$

\mbox{定理~29.4.3} 的证明中最困难的一步，是下面的引理。

\textbf{29.4.5. \trhLemma。}\ ——\ \emph{设 $\pi\colon X\to Y$ 是余维数为 $r$ 的纯维\trProjective{} $k$-\trScheme{}之间的一个\trClosedEmbedding，且 $Y$ 满足对 $\mathrm{Ext}$ 的函子性\trSerreDuality。那么，我们有随 $\mathscr{F}\in\mathit{Coh}_X$ 函子的同构}
\begin{equation*}
\tag{29.4.5.1}\label{eq:hom-ext-iso-lemma}
\mathrm{Hom}_X(\mathscr{F},\pi^{-1}\mathscr{E}xt^r_Y(\pi_*\OO_X,\omega_Y))\overset{\sim}{\longleftrightarrow}\mathrm{Ext}^r_Y(\pi_*\mathscr{F},\omega_Y).
\end{equation*}

\emph{证明。}\quad 选取 $\omega_Y$ 的一个内射分解
\[
0\longrightarrow\omega_Y\longrightarrow\mathscr{I}^0\longrightarrow\mathscr{I}^1\longrightarrow\mathscr{I}^2\longrightarrow\cdots
\]
为了计算 \eqref{eq:hom-ext-iso-lemma} 的右边，我们从这一分解中去掉 $\omega_Y$，然后施加 $\mathrm{Hom}_Y(\pi_*\mathscr{F},\cdot)$，便得到
\[
0\longrightarrow\mathrm{Hom}_Y(\pi_*\mathscr{F},\mathscr{I}^0)\longrightarrow\mathrm{Hom}_Y(\pi_*\mathscr{F},\mathscr{I}^1)\longrightarrow\mathrm{Hom}_Y(\pi_*\mathscr{F},\mathscr{I}^2)\longrightarrow\cdots,
\]
或者（由 $\pi$ 为\trClosedEmbedding{}时 $\pi_*$ 与 $\pi^!$ 在 $\OO$-\trModule{}上的\trAdjointness，\mbox{练习~29.4.B(a)}）：
\[
0\longrightarrow\mathrm{Hom}_X(\mathscr{F},\pi^!\mathscr{I}^0)\longrightarrow\mathrm{Hom}_X(\mathscr{F},\pi^!\mathscr{I}^1)\longrightarrow\mathrm{Hom}_X(\mathscr{F},\pi^!\mathscr{I}^2)\longrightarrow\cdots.
\]
受此启发，考虑复形
\begin{equation*}
\tag{29.4.5.2}\label{eq:pi-bang-complex}
0\longrightarrow\pi^!\mathscr{I}^0\longrightarrow\pi^!\mathscr{I}^1\longrightarrow\pi^!\mathscr{I}^2\longrightarrow\cdots.
\end{equation*}
首先请注意，\eqref{eq:pi-bang-complex} 确实是一个复形；其次请注意，诸 $\pi^!\mathscr{I}^i$ 都是内射的 $\OO_X$-\trModule（由\mbox{练习~29.4.B(b)}）。

\textbf{29.4.D.（练习）}\quad 证明：\eqref{eq:pi-bang-complex} 在第 $i$ 步处的\trCohomology{}是 $\pi^{-1}\mathscr{E}xt^i_Y(\pi_*\OO_X,\omega_Y)$。

于是，由\mbox{命题~29.4.4}，复形 \eqref{eq:pi-bang-complex} 在第 $r$ 步之前是正合的。

\textbf{29.4.E.（练习）}\quad 证明：存在一个直和分解 $\pi^!\mathscr{I}^r=\mathscr{J}\oplus\mathscr{K}$，使得
\[
0\longrightarrow\pi^!\mathscr{I}^0\longrightarrow\pi^!\mathscr{I}^1\longrightarrow\cdots\longrightarrow\pi^!\mathscr{I}^{r-1}\longrightarrow\mathscr{J}\longrightarrow 0
\]
正合。提示：先做出 $r=1$ 的情形。再一条提示：注意，若
\[
0\longrightarrow\mathscr{I}^0\overset{\alpha}{\longrightarrow}\mathscr{F}
\]
正合，而 $\mathscr{I}^0$ 是内射的，则存在一个「分裂」$\alpha$ 的映射 $\beta\colon\mathscr{F}\to\mathscr{I}^0$，它使我们能把 $\mathscr{F}$ 写成直和 $\mathscr{I}^0\oplus\mathscr{J}^1$。若 $\mathscr{F}$ 进而还是内射的，则 $\mathscr{J}^1$ 也是内射的。这便是一个归纳法的开端。

\textbf{29.4.F.（练习）}\quad （我们沿用上一道练习的记号。）请把 $\ker(\mathscr{K}\to\pi^!\mathscr{I}^{r+1})$ 识别为 $\pi^{-1}\mathscr{E}xt^r_Y(\pi_*\OO_X,\omega_Y)$。

\textbf{29.4.G.（练习）}\quad 把上面的各块拼合起来，以完成\mbox{引理~29.4.5} 的证明。\hfill$\square$

我们现在已经为证明\mbox{定理~29.4.3} 做好了准备。

\textbf{29.4.6.}\ \emph{\mbox{定理~29.4.3} 的证明。}\quad 设 $\mathscr{F}$ 是 $X$ 上的一个\trCoherentSheaf。我们希望描述一个同构
\[
\mathrm{Hom}_X(\mathscr{F},\pi^{-1}\mathscr{E}xt^r_Y(\pi_*\OO_X,\omega_Y))\overset{\sim}{\longleftrightarrow}\mathrm{H}^n(X,\mathscr{F})^{\vee},
\]
它随 $\mathscr{F}$ 函子。由\mbox{引理~29.4.5}，我们有一个函子性的同构
\[
\mathrm{Hom}_X(\mathscr{F},\pi^{-1}\mathscr{E}xt^r_Y(\pi_*\OO_X,\omega_Y))\overset{\sim}{\longleftrightarrow}\mathrm{Ext}^r_Y(\pi_*\mathscr{F},\omega_Y).
\]
由 $Y$ 上对 $\mathrm{Ext}$ 的函子性\trSerreDuality，我们有一个函子性的同构
\[
\mathrm{Ext}^r_Y(\pi_*\mathscr{F},\omega_Y)\overset{\sim}{\longleftrightarrow}\mathrm{H}^{N-r}(Y,\pi_*\mathscr{F})^{\vee}.
\]
由于 $\mathscr{F}$ 是 $X$ 上的\trCoherentSheaf，而 $N-r=n$，右边正是 $\mathrm{H}^n(X,\mathscr{F})^{\vee}$。\hfill$\square$

\textbf{29.4.7. \mbox{定理~29.4.3} 的应用。}\quad 我们首先把\mbox{定理~29.4.3} 应用于 $X$ 是 $Y$ 上的一个\trEffectiveCartierDivisor{}这一特殊情形。我们可以这样计算\trDualizingSheaf{} $\pi^{-1}\mathscr{E}xt^1_Y(\pi_*\OO_X,\omega_Y)$：用 $\pi_*\OO_X$ 的任意一个局部自由分解（在 $Y$ 上）来计算 $\mathscr{E}xt^1_Y(\pi_*\OO_X,\omega_Y)$（\mbox{练习~29.2.E}）。但 $\pi_*\OO_X$ 有一个特别简单的分解，即 $X$ 的\trClosedSubscheme{}正合列~\mbox{(9.1.2.1)}：
\begin{equation*}
\tag{29.4.7.1}\label{eq:cartier-closed-subscheme-seq}
0\longrightarrow\OO_Y(-X)\longrightarrow\OO_Y\longrightarrow\pi_*\OO_X\longrightarrow 0.
\end{equation*}
我们这样计算 $\mathscr{E}xt^{\bullet}(\pi_*\OO_X,\omega_Y)$：截断（去掉 $\pi_*\OO_X$），然后施加 $\mathscr{H}om(\cdot,\omega_Y)$。$\mathscr{E}xt^{\bullet}(\pi_*\OO_X,\omega_Y)$ 是如下复形的上同调
\[
0\longrightarrow\mathscr{H}om_Y(\OO_Y,\omega_Y)\longrightarrow\mathscr{H}om(\OO_Y(-X),\omega_Y)\longrightarrow 0,
\]
也就是说，是
\[
0\longrightarrow\omega_Y\longrightarrow\omega_Y\otimes\OO_Y(X)\longrightarrow 0.
\]
的上同调。

我们立即看到：当 $i\ne 0$、$1$ 时，$\mathscr{E}xt^i(\pi_*\OO_X,\omega_Y)=0$。此外，由\mbox{命题~29.4.4}，$\mathscr{E}xt^0(\pi_*\OO_X,\omega_Y)=0$。（一条不重要的题外话：你可以用它来证明 $\omega_Y$ 没有\trEmbeddedPoint。）

我们现在考虑 $\omega_X=\coker(\omega_Y\to\omega_Y\otimes\OO_Y(X))$。用\trInvertibleSheaf{} $\OO_Y(X)$ 张量 \eqref{eq:cartier-closed-subscheme-seq}，随后再与 $\omega_Y$ 作张量，便得到
\[
\omega_Y\longrightarrow\omega_Y\otimes\OO_Y(X)\longrightarrow\omega_Y\otimes\pi_*(\OO_Y(X)|_X)\longrightarrow 0.
\]
右边的项 $\omega_Y\otimes\pi_*(\OO_Y(X)|_X)$ 常常（多少不太正式地）被写作 $\omega_Y(X)|_X$。于是
\[
\omega_X=\coker(\omega_Y\to\omega_Y\otimes\OO_Y(X))=\omega_Y(X)|_X,
\]
而且这一识别是\emph{\trCanonical{}的}。

我们便证明了下述结果。

\textbf{29.4.8. \trhProposition（\trAdjunctionFormula）。}\ ——\ \emph{设 $Y$ 是\trPureDimension{}为 $n$ 的 \trCM{} \trProjectiveScheme（由\mbox{推论~29.3.14}，它满足对 $\mathrm{Ext}$ 的函子性\trSerreDuality），而 $X$ 是 $Y$ 上的一个\trEffectiveCartierDivisor。（于是由 \S26.2.4，$X$ 是\trCM{}的，从而由\mbox{推论~29.3.14}，它满足对 $\mathrm{Ext}$ 的函子性\trSerreDuality。）则 $\omega_X=\omega_Y(X)|_X$。特别地，若 $\omega_Y$ 是 $Y$ 上的一个\trInvertibleSheaf，则 $\omega_X$ 是 $X$ 上的一个\trInvertibleSheaf。}

作为一个直接的应用，我们有下面的结果。

\textbf{29.4.H.（练习）}\quad 设 $X$ 是 $\mathbb{P}^n$ 中的一个\trCompleteIntersection，由次数为 $d_1$，…，$d_r$ 的诸\trHypersurface{}给出。（注意：由\mbox{命题~26.2.6}，$X$ 是\trCM{}的，从而由\mbox{推论~29.3.14}，它满足对 $\mathrm{Ext}$ 的函子性\trSerreDuality。）证明：$\omega_X\cong\OO_X(-n-1+\sum d_i)$。进而，若 $X$ 还是\trSmooth{}的，证明：$\omega_X\cong\det\Omega_X$。（最后一句话的提示：使用 $\mathscr{K}$ 的\trAdjunctionFormula，即\mbox{练习~21.5.B}。）

但我们还能说得更多。

\textbf{29.4.I.（练习）}\quad 设 $\pi\colon X\hookrightarrow Y$ 是余维数为 $r$ 的一个\trRegularEmbedding，其正规层为 $\mathscr{N}_{X/Y}$（由\mbox{命题~21.2.16(b)}，它是\trLocallyFree{}的）。设 $\mathscr{L}$ 是 $Y$ 上的一个\trInvertibleSheaf。

(a) 证明：当 $i\ne r$ 时，$\mathscr{E}xt^i_Y(\pi_*\OO_X,\mathscr{L})=0$。

(b) 请描述一个\trCanonical{}同构
\[
\pi^{-1}\mathscr{E}xt^r_Y(\pi_*\OO_X,\mathscr{L})\overset{\sim}{\longleftrightarrow}(\det\mathscr{N}_{X/Y})\otimes_{\OO_X}\mathscr{L}|_X.
\]

两部分共同的提示：先处理 $\mathscr{L}=\OO$ 的情形，然后注意到：假如你已经细心而正确地解出了 (b)，这个问题便是局部的。

（注意：由于\mbox{练习~29.4.I} 与对偶性并无显式关联，我们对 $Y$ 不作任何\trProjective{}性假设；这完全是一个局部的问题。）由\mbox{练习~29.4.I}，我们推出下面的结论。

\textbf{29.4.J.（重要练习）}\quad 设 $X$ 是 $\mathbb{P}^n_k$ 中余维数为 $r$ 的一个\trRegularEmbedding。（此时 $X$ 由\mbox{命题~26.2.6} 是\trCM{}的，从而由\mbox{推论~29.3.14} 满足对 $\mathrm{Ext}$ 的函子性\trSerreDuality。）证明：
\[
\omega_X=\pi^{-1}\mathscr{E}xt^r_{\mathbb{P}^n}(\pi_*\OO_X,\omega_{\mathbb{P}^n})\cong(\det\mathscr{N}_{X/\mathbb{P}^n})\otimes\omega_{\mathbb{P}^n}|_X.
\]
特别地，$\omega_X$ 是一个\trInvertibleSheaf。

\textbf{29.4.9.}\ \emph{题外话：Gorenstein 奇点。}\quad 若纯维的 \trCM{} \trProjective{} $k$-\trScheme{} $X$ 的 $\omega_{X/k}$ 是一个\trInvertibleSheaf，我们就说 $X$ 是 \emph{Gorenstein} 的。借助\mbox{命题~29.4.8}，你可以证明：这是一个可以在纯维\trProjective{} $k$-\trScheme{}的各个\trClosedPoint{}处按茎局部检验的条件，正如 \trCM{}性的定义那样。你能写出这样一个定义吗？在正确的普遍性下，Gorenstein 奇点的定义当然并不要求所论的\trScheme{}是射影的。这段题外话将引导你借助一个「\trSlicing{}判别法」，得到一个正确却谈不上优雅的定义。一切「局部\trCompleteIntersection」的 $k$-\trScheme{}都是 Gorenstein 的，而一切 Gorenstein 的 $k$-\trScheme{}按定义都是\trCM{}的。关于 Gorenstein 环与 Gorenstein 概形的一份极佳入门材料，见 [BruH, §3]。

\textbf{29.4.K.（重要练习）}\quad 设 $X$ 是 $\mathbb{P}^n_k$ 的一个\trSmooth{}且纯余维数为 $r$ 的\trSubvariety（因而由\mbox{练习~13.2.M(b)}，它是一个\trRegularEmbedding）。证明：$\omega_X\cong\mathscr{K}_X$。提示：两边都满足\trAdjunctionFormula（关于 $\Omega$ 的\trAdjunctionFormula，见\mbox{练习~21.5.B}）：它们同构于 $(\det\mathscr{N}_{X/Y})\otimes(\omega_{\mathbb{P}^n}|_X)\cong(\det\mathscr{N}_{X/Y})\otimes(\mathscr{K}_{\mathbb{P}^n}|_X)$。\footnote{译注：原书这两处下标作「$X/Y$」，其中 $Y$ 未在本练习中出现（本练习中 $X\subset\mathbb{P}^n_k$），疑当作「$X/\mathbb{P}^n$」，兹照录原书。——译者注}

\textbf{29.4.10.}\ \emph{\trhRemark{}。}\quad 正如久已许诺的那样：把\mbox{推论~29.3.10} 与\mbox{练习~29.4.K} 结合起来，\mbox{定理~18.5.1} 与\mbox{进一步的奇迹~18.5.2} 中给出的那个版本的\trSerreDuality{}便随之得证。


\epigraph{在那些乐于为理解而忧心、追逐完美的人眼中，一部作品永远不会完成——「完成」一词于他们毫无意义——而只是被放弃。而这种放弃——正是它把作品交付给烈火或公众（或许是出于倦怠，或许是出于交稿期限）——在他们看来不过是某种意外罢了，恰似叩门声骤然响起，惊破了他们的梦境……}{——保罗·瓦莱里，[Coh, pp.~7--8]（Tadashi Tokieda 英译）}

\epigraph{好了，我已经给你们讲了两年课，现在就要收场了。在某些方面我想道个歉，在另一些方面则不想。我希望——事实上，我知道——你们当中能有二三十人一直满怀兴奋地听懂了全部内容，并从中获得了乐趣。但我也明白：「讲授的效力微乎其微，除非身处那种讲授几乎纯属多余的幸福境遇。」所以，对那二三十位听懂了全部内容的人，请允许我说一句：我其实什么也没做，不过是把那些东西展示给了你们。至于其余的人，倘若我让你们恨上了这门学科，我很抱歉。我以前从未教过基础物理，在此致歉。我只希望自己没有给你们造成严重的困扰，也希望你们不要就此离开这件令人兴奋的事业。我希望将来能有别人以一种不让你们消化不良的方式把这门课教给你们，而有朝一日你们终究会发现：它并不像看上去那么可怕。}{——R.~费曼，[FLS, Epilogue]}

